\documentclass[11pt]{article}
\usepackage{booktabs}
\usepackage[hidelinks]{hyperref}
\title{Checking warehouse movement before crediting recovery}
\date{}
\begin{document}
\maketitle

Distinct occupied cells at a tick boundary do not establish
one-cell-at-a-time motion. On the
\href{https://github.com/eternal-november/eos-research-programmes/pull/402}
{existing development cases}, the native simulator applied
moves that violate that constraint. The advantage of ordinary
replanning over waiting survives in aggregate when both arms
are given the same stricter movement rule. Most runs still do
not meet the full service deadline.

\section{What was held fixed}
The public MRWS source is pinned at
\href{https://github.com/automaticdai/MRWS_MultiRobot_Warehouse_Scheduling/tree/941338e9ea4c4eaebb87f8374f80fd576497ad01}
{\texttt{941338e9}}. The replay imports the previous frozen
harness: same sixteen development seeds, source layout, order
offers, per-actor and per-tick fault opportunities, and
terminal, waiting and replanning policies. The proposed change
is confined to a reviewed motion check in the local simulator.
The verifier observes individual move calls and does not add
new information to the robot or to the native planner. The
replay matches offered order identities, content, priority and release times
on all continuing arms, as well as fault opportunity tables,
not actual fault claims that can change with policy.

The native occupied-cell and wall checks do not test adjacency.
In the source a refuge move can finish with a path still cached
from the previous cell; a subsequent action can consume that
path without checking its origin again. This explains a source
mechanism, not the cause of every non-unit move recorded below.
The audit counts applied position updates relative to the
immediately preceding position of that robot, not its last
tick-boundary location. It separately counts repeated applied
moves by the same robot in one clock tick.

In addition to the untouched native motion rule, both controller
arms run under two stricter laws. The unit-step law admits only
Manhattan distance exactly one per accepted position update.
The unit-per-tick law additionally admits at most one accepted
move per robot per tick, even if a different actor's deadlock
routine would have made the move. These are explicit alternate
plant assumptions, not corrected descriptions of what the
upstream simulator already guarantees. A motion refusal leaves
position and carried inventory unchanged; it clears a stale
cached path and ends the current action. Earlier moves remain
in effect. The native pinned plant remains available unchanged.

\section{Movement and service results}

\begin{center}
\begin{tabular}{lrrr}
\toprule
Accepted movement rule & Native stop policy & Wait & Replan\\
\midrule
Unchanged native & 2 & 34 & 49\\
Unit step per move & 4 & 35 & 49\\
One unit step per robot per tick & 3 & 34 & 48\\
\bottomrule
\end{tabular}

Counts are completed orders across sixteen paired development runs
per cell, not percentages or a guarantee.
\end{center}

The unchanged native motion arm accepted 69 non-unit updates
with waiting, and 71 with replanning; 16 and 15 runs,
respectively, contained at least one. The stricter arms accepted
none. The unit-step-only rule still applied 38 extra same-robot
moves within a tick with waiting and 41 with replanning. The
stronger rule applied none. The last row had 14 timed-out
waiting runs and 10 timed-out replanning runs. In particular,
preserving distinct occupied cells is not sufficient to transport
a service number to a plant that limits individual motion.
The original aggregate direction survives both checks on these
seeds, without establishing dominance or field fidelity.

The replay also tracks calls to the native ordinary motion routine
that end without an applied move after querying a cell made
unavailable by the cached fault neighborhood alone, not by a
wall or current occupancy. There were 1,856 such calls in the
native-motion replanning arm. This count is not a causal allocation
of timeouts: it is a trace indicator, and an obstacle snapshot
can differ from the state after other moves in a tick. It shows
why an occupied-destination rejection is not the only condition
under which the native planner can stop making progress.

This adds a motion-fidelity qualification and a direct baseline
sensitivity analysis. No independent source selection, I1 acceptance,
or proof of superiority over a matched planner is claimed.
\end{document}

